\scanpage{278}
\mspage{8}

et que $\operatorname{Loc}(2^o)$ est connu par $\operatorname{Loc}(1,2,3)\Longleftrightarrow\operatorname{Loc}(1,2^o,3^o)$. Il reste donc à prouver que si

\[
f:X\unsure{{-}|\longrightarrow}Y\qquad\text{dans }\mathrm{Cat}_{/S}
\]

est telle que \deleted{$\unsure{{}_{s}f}$}

\[
{}_{s\backslash}f:{}_{s\backslash}X{-}{}_{s\backslash}Y
\]

est dans $W$ pour tt $s\in\operatorname{Ob}S$, alors $f\in W$. Considérons \unsure{alors|donc} le carré commutatif

\diagram{\begin{tikzcd}[ampersand replacement=\&]
X \arrow[r,"{\unsure{i_X|\iota_X}}"] \arrow[d,leftrightarrow,"f"'] \& \Phi_S(X) \arrow[d,"{\Phi_S(f)}"]\\
Y \arrow[r,"i_Y"] \& \Phi_S(Y),
\end{tikzcd}}

où $i_X,i_Y\in W$. par $\operatorname{Loc}(3'\,\mathrm{bis}\ b)$. L'hypothèse sur $f$ signifie que $\Phi_S(f)$ induit des $W$-isomorphismes sur les fibres, donc par $\operatorname{Loc}(3'\,\mathrm{bis})$ $\Phi_S(f)\in W$, donc $f\in W$ par $\operatorname{Loc}(1)$, qed.

\uline{2. Axiomes} \uline{Loc} (4) \uline{et} \uline{Loc} (5) (des sommes \unsure{directes|discrètes|disjointes}, de \unsure{connexité?})

% slanted
\margin{Il s'agit ici d'axiomes qu'on peut considérer comme étant de nature plus ou moins ``\unsure{triviale}''. Dans \unsure{tous les} cas à ma connaissance, ils sont bel et bien trivialement vérifiés.}

\uline{Loc 4} \inserted{(Axiome des sommes \unsure{discrètes|directes})} Pour une famille $(f_i)_{i\in I}$ de flèches de Cat, \unsure{si} $f=\coprod f_i$, on a $f\in W$ sss $f_i\in W$ pour tt $i\in I$.

Cet axiome se décompose en deux, en le ``\unsure{Si|si}'' et le ``seulement si'', que je vais \deleted{\unsure{désigner}} formuler sous la forme la plus économique :

\uline{Proposition} 2.\unsure{1|4}. L'axiome \uline{Loc} 4 \inserted{des sommes} \unsure{revient|équivaut} à la conjonction des deux axiomes suivants :

% archive sheet number: 277
